 \subsection{Minimal Explantation Argument}
 \label{subsection:explantion}
We are about to prove the main lemma stating that the number of faults causing an occurrence of slice is linear in both the slice's size and the minimal number of edges required to span a Toom's contour that induces it. First, we define the explanation tree of a slice to be the subtree of its history that connects its poles. The relation between the leaves of the subtree and its arrows, forks, and the size of its origin slice is embedded as a requirement in the definition. Then, we prove that, assuming no step admits an infinite slice, any slice has an explanation tree.
  \begin{definition}
    \label{def:explantation}
    Consider a Tooms contour graph induced by Toric encoded circuit, with sweep-cycle parameter $\xi = 1$, and  let $\alpha, \beta$ be constants and $\mathcal{X}_{0}, \mathcal{X}_{1}, \mathcal{X}_{2}, \mathcal{X}_{3} \in \mathbb{R}$ be quadruple, defining such:
    \begin{equation*}
      \begin{split}
      \alpha &\colonequals 40(d+1)^{2}, \ \ \  \beta \colonequals 5 (d+1) \\
      \mathcal{X}_{r} &\colonequals 1 - r \frac{d+1}{\alpha} = 1 -  \frac{r}{40(d+1)}
      \end{split}
    \end{equation*}
In addition, we would like to refer to $\mathcal{X}_{0}, \mathcal{X}_{1}, \mathcal{X}_{2}, \mathcal{X}_{3}$ in a cyclic manner, so we define:
    \begin{equation*}
      \begin{split}
        \mathcal{X}_{-1} \colonequals \mathcal{X}_{3}
      \end{split}
    \end{equation*}
Given sets $W$ of points, $A$ of arrows, and $F$ of forks, we define $U$ as the set of vertices of the graph induced by the edges of $A \cup F$. The sets $W$, $A$, and $F$ form an explanation of a slice $K$, with poles $v_{1}, \ldots, v_{d+1}$, and order $\text{order}(K) = r$, if and only if:
    \begin{enumerate}
      \item $W \cup \{v_{1} ,.., v_{d+1}\} \subset U \subset \History{K}$ and the graph $(U, A\cup F)$ is connected. 
      \item All elements of $W$ are leaves. 
      \item For $m = |W|$ and $s=\Size{K}$ we have $|A| \le \alpha (m - \mathcal{X}_{r}) - \beta s $ and $|F| < m -1 $. 
    \end{enumerate}
Furthermore, for a step $\tau$ and a slice $K$ in the analysis Toom's contour at step $\tau$, we extend the definition of explanation as above, but require the following weaker inequality instead: For $m = |W|$ and $s = \Size{K}$, we have $|A| \le \alpha (m + 1) - \beta s$ and $|F| < m - 1$.
  \end{definition}
  \begin{claim}
    \label{claim:edgesin}
Let $P_{1}, \ldots, P_{\mu+1}$ be the reduced Pauli-path that induces Toom's contour. And let $t \le \mu+1$ such that $P_{1}, \ldots, P_{t}$ do not admit an infinite slice. Every slice $K$, with time coordinate $\Time{K} \le t$, has an explanation. Moreover, if a step $\tau+1$ is the first step to admit an infinite slice, any slice in the analysis of Toom's contour up to step $\tau$ has an explanation.
  \end{claim}

  \begin{proof}
Let $h = \Time{K} - \min_{u \in \History{K}} \Time{u}$. For slices in a standard Toom's contour, we prove by induction on $h$.
    \begin{enumerate}
      \item Base: $h=0$, meaning $T$ is constant on $\History{K}$, so no arrows connect points of $\History{K}$. Thus, $\History{K}$, as well as $\mathbf{K}$, consists only of a single leaf, namely $K$ is a result of a fault, and $r = \text{order}(\mathcal{X}_{r}) \in \{0,2\}$. In that case, there is a single vertex $u \in V_{\mathcal{H}}$ such that $U = W = \{u\}$, so $m=1$, $s=0$, and $|A|, |F| = 0$. It is easy to check that those values satisfy the requirement. 
      \item Induction Step: $h \neq 0$ implies $K \neq \History{K}$. Combining with the fact that no infinite slice is admitted, by \Cref{claim:prevslices}, we know that for some $k$ there exist slices $K_{0}, \ldots, K_{k}$ and forks $F_{1}, \ldots, F_{k}$ that satisfy \Cref{claim:prevslices}. In particular, for $r < 3$ it holds that:
        \begin{equation*}
          \begin{split}
            \sum_{i=0}^{k}\Size{K_{i}} +  \sum^{k}_{i=1}\Size{F_{i}} \ge \Size{K}  
          \end{split}
        \end{equation*}
        Whereas, for $r=3$: 
        \begin{equation*}
          \begin{split}
            \sum_{i=0}^{k}\Size{K_{i}} +  \sum^{k}_{i=1}\Size{F_{i}} \ge \Size{K} + \xi = \Size{K} + 1 
          \end{split}
        \end{equation*}
        Consider $i \in \{0, \ldots, k\}$, and denote $s_{i} = \Size{K_{i}}$. By the inductive hypothesis and since $\Time{K_{i}} = \Time{K} - 1$, we know that $K_{i}$ has an explanation formed by a set of leaves $\mathbf{W}_{i}$, a set of arrows $\mathbf{A}_{i}$, and a set of forks $\mathbf{F}_{i}$. Let $m_{i} = |\mathbf{W}_{i}|$. Thus, the sets $\mathbf{A}_{i}$ and $\mathbf{F}_{i}$ have at most $\alpha (m_{i} - \mathcal{X}_{r-1}) - \beta s_{i}$ arrows and $m_{i} - 1$ forks. We define:
        \begin{enumerate}
          \item $\mathbf{W} = \cup_{i=0}^{k} \mathbf{W}_{i}$, Notice that $\mathbf{W}_{i} \subset K_{i}$ and that $\Time{K_{i}} = \Time{K_{j}}$ for all $i,j$, thus by \Cref{claim:disj} $\mathbf{W}$ is a union of disjoint set, Hence $W$ contains  $m = \sum_{i=0}^{k}m_{i}$ leaves. 
          \item $\mathbf{F} = \{ F_{1},..,F_{k} \} \cup_{i = 0}^{k}\mathbf{F}_{i}$ with at most $ k + \sum_{i=0}^{k}{m_{i} -1 } = m- 1$ elements. 
          \item $\mathbf{A} = \{ \{v_{i}, \prei{v_{i}} \} \} \cup_{i=0}^{k}{ \mathbf{A}_{i} }$. For bounding the number of arrows in $\mathbf{A}$ we condition upon whether $r = 3$, For convenient, let us denote by $Y$ the indicator, which indicate if $r = 3$.   
            \begin{equation*}
              \begin{split}
                & d+1 + \sum^{k}_{i=0}{\alpha (m_{i}-\mathcal{X}_{r-1})  - \beta s_{i} } \\ 
                \le & d+1 + \alpha m - \mathcal{X}_{r-1}\alpha(k+1) - \beta\left(s + Y\xi - k \cdot \Size{\text{fork}}  \right) \\
                = & \alpha \left(m - \mathcal{X}_{r}\right) - \beta s + \left(\alpha ( \mathcal{X}_{r} - \mathcal{X}_{r-1} ) + d+1 -  \mathcal{X}_{r-1} k \alpha - \beta Y \xi + \beta k \cdot \Size{\text{fork}} \right)
              \end{split}
            \end{equation*}
So it's left to show that the term in the closure is negative for all integers $k$:
            \begin{equation*}
              \begin{split}
\alpha ( \mathcal{X}_{r} - \mathcal{X}_{r-1} ) + d+1 -  \mathcal{X}_{r-1} k \alpha - \beta Y \xi + \beta k \cdot \Size{\text{fork}} 
              \end{split}
            \end{equation*}
            We do that by showing that for the choice of the parameter values, $\alpha, \beta, \mathcal{X}_{0}, \ldots, \mathcal{X}_{3}$, the slope is negative, and then the value of the term at $k$ equals zero is negative.  For the slope to be negative, we require: 
            \begin{equation*}
              \begin{split}
                - \alpha \mathcal{X}_{r-1}  +  \beta \Size{\text{fork}} \le 0 & \Rightarrow   \Size{\text{fork}} \le \frac{\alpha}{\beta} \mathcal{X}_{r-1} 
              \end{split}
            \end{equation*}
            For negativity at the origin, $k=0$, we require for $r < 3$:
            \begin{equation*}
              \begin{split}
                (\mathcal{X}_{r}-\mathcal{X}_{r-1})  & \alpha  + d + 1  \le 0 
                \\ &  \Rightarrow \mathcal{X}_{r} \le  \mathcal{X}_{r-1} - \frac{d+1}{\alpha}  
  \end{split}
\end{equation*}
And for $r =3$:
\begin{equation*}
  \begin{split}
(\mathcal{X}_{3}-\mathcal{X}_{2}) & \alpha  + d + 1  -\beta \xi \le 0  \\ 
& \Rightarrow \mathcal{X}_{3} \le  \mathcal{X}_{0} - 3\frac{d+1}{\alpha} + \frac{\beta}{\alpha}\xi  
              \end{split}
            \end{equation*}
So in total: 
\begin{equation*}
  \begin{split}
&  \mathcal{X}_{1} \le  \mathcal{X}_{0} - \frac{d+1}{\alpha}  \\
&  \mathcal{X}_{2} \le  \mathcal{X}_{0} - 2\frac{d+1}{\alpha}  \\ 
&  \mathcal{X}_{3} \le  \mathcal{X}_{0} - 3\frac{d+1}{\alpha} + \frac{\beta}{\alpha}\xi  \\
&  \mathcal{X}_{0} \le  \mathcal{X}_{0} - 4\frac{d+1}{\alpha}  + \frac{\beta}{\alpha}\xi  \\
 & \Size{\text{fork}} \le \frac{\alpha}{\beta} \mathcal{X}_{r} 
  \end{split}
\end{equation*}
The first three inequalities hold by the definitions of the $\mathcal{X}_{r}$, the fourth is satisfied since $\frac{\beta}{\alpha}\xi = 5 \frac{d+1}{\alpha}$, the last is obtained due to \Cref{claim:maxfork}, that states that the maximal size of a fork is at most $2d$, and that $\frac{\alpha}{\beta}\mathcal{X}_{r}$ is greater than $4d$.
        \end{enumerate}
         Finally, by \Cref{claim:prevslices}, the graphs induced by $\mathbf{A}_{i} \cup \mathbf{F}_{i}$ are connected, thus the graph induced by $\mathbf{A}\cup \mathbf{F}$ is also connected. Therefore, $\mathbf{W}, \mathbf{A}, \mathbf{F}$ form an explanation of $K$. 
    \end{enumerate}
It's left to prove for slices in the analysis of Toom's contours. Notice that the 'forming rules' of the analysis of Toom's contour are different from the standard Toom's contour only at the last layer. Therefore, the induction proof captures all the slices not in the last layer. For any slice at the last layer, we repeat the induction step while using \Cref{claim:shrinkan} instead of \Cref{claim:prevslices}, so the bounding over $\mathbf{A}$'s size becomes:
    \begin{equation*}
      \begin{split}
        |\mathbf{A}|  & \le \alpha \left(m + 1 \right) - \beta s \\ 
        & + \left( - \alpha (  1 +  \mathcal{X}_{r-1} ) + d+1 -  \mathcal{X}_{r-1} k \alpha + d + d^{\prime} + 1 + \beta k \cdot \Size{\text{fork}} \right)  \\ 
        & \le \alpha \left(m + 1 \right) - \beta s \\ 
        & + \left( - \alpha  + d+1 -  \mathcal{X}_{r-1} k \alpha + d + d^{\prime} + 1 + \beta k \cdot \Size{\text{fork}} \right)  
      \end{split}
    \end{equation*}
    Where we used the fact that $\mathcal{X}_{r-1} \in (0,1)$ for $r \in \{0,1,2,3\}$. Clearly, the bound is satisfied since $\alpha \ge 2d + d^{\prime} + 1$ and $\mathcal{X}_{r-1} \alpha \ge \beta \cdot \Size{\text{fork}}$.
  \end{proof}
We summarize by using the linear relation between the leaves of a Toom's contour and the faults that induce it. That gives an upper bound on the probability of a slice occurring when the noise is local stochastic.
  \begin{corollary}
    \label{coro:edgefaultratio}
Let $K$ be a slice with size $s$, and explanation with $|A|$ arrows, and $|F|$ forks. Then it is induced by at least the following number of faults:
    \begin{equation*}
      \begin{split}
        \frac{1}{2^{\max\{d^{\prime}, d - d^{\prime}\}}}  \frac{1}{2} \left( |F|+1 +  \alpha^{-1} ( |A|  + \beta s - 1)  \right)
      \end{split}
    \end{equation*}
    Furthermore, the ratio of leaves to edges in an explanation is greater than $\frac{1}{2\alpha} = \frac{1}{80(d+1)^{2}}$.
  \end{corollary}
  \begin{proof}
We bound \Cref{claim:edgesin} by the number of leaves in the explanation and use the fact that any $d^{\prime}$-face has $2^{d^{\prime}}$ vertices, so each $X$-fault and $Z$-fault contributes at most $2^{d^{\prime}}$ and $2^{d-d^{\prime}}$ leaves to the explanation.
  \end{proof}

    %\ctt{ Check what happened if we set $\xi = \sqrt{d}, \beta = \sqrt{d}$ is that better? }


